\section*{Chapitre \ref{ch:rs:market-data-for-research} --- Données de marché pour la recherche}

\begin{solution}{exo:rs:market-data-for-research:1}
$(300 \times 20.00 + 500 \times 20.02 + 200 \times 19.99)/1\,000 = 20\,008/1\,000 = 20.008$.
\end{solution}

\begin{solution}{exo:rs:market-data-for-research:2}
La fourchette vaut $0.01/20 = 0.0005$ en termes logarithmiques, de sorte que chaque rendement d'une transaction à la suivante gagne $0.0005^2/2 = 1.25 \times 10^{-7}$ de variance~;
sur 20\,000 rendements, cela donne $0.0025$. Avec la vraie valeur $0.02^2 = 0.0004$, la somme vaut $0.0029$~: une volatilité annoncée de 5.39~\%, 2.7
fois la vraie valeur.
\end{solution}

\begin{solution}{exo:rs:market-data-for-research:3}
$40 \times 1.5 \times 10^6 \times 5\,000 \times 252 = 7.56 \times 10^{13}$ octets, soit environ 76 téraoctets par an avant compression.
\end{solution}

\begin{solution}{exo:rs:market-data-for-research:4}
$3(1 + 0.8^2) = 4.92$. Un kurtosis de 6 exige $1 + \mathrm{CV}^2 = 2$, soit un coefficient de variation de 1.
\end{solution}

\begin{solution}{exo:rs:market-data-for-research:5}
Position détenue~: 70, 71, 76 (le contrat suivant à partir du troisième jour). L'écart de roll au deuxième jour vaut $74 - 71 = 3$. Ajusté rétroactivement~: 73, 74, 76.
Série ajustée proportionnellement~: $70 \times 74/71 = 72.96$, $74$, $76$. Le rendement de la position détenue au troisième jour est $76/74 - 1 = 2.7~\%$, que la
série ajustée proportionnellement montre et que la série brute du contrat le plus proche ($72/71$ puis un saut) ne montre pas.
\end{solution}

\begin{solution}{exo:rs:market-data-for-research:6}
Une même valeur échangée correspond à dix fois plus d'actions à USD~10, si bien qu'un seuil exprimé en actions clôt dix fois plus de \omterm{def:rs:market-data-for-research:bar}{barres} par jour en
2015 qu'en 2025. Les \omterm{def:rs:market-data-for-research:bar}{barres de montant} clôturent le même nombre dans les deux années~: c'est la valeur, et non le nombre d'actions, qui mesure le volume des échanges.
\end{solution}

\begin{solution}{exo:rs:market-data-for-research:7}
3\,472 \omterm{def:rs:market-data-for-research:bar}{barres} de 16.0 transactions en moyenne (la \omterm{def:rs:market-data-for-research:bar}{barre} médiane en compte 7), avec un kurtosis de 35.4 pour les rendements de mid (24.2 de clôture à clôture).
Le seuil est réestimé à partir des \omterm{def:rs:market-data-for-research:bar}{barres} qui viennent de se clore. Parce que les signes des ordres sont autocorrélés (métaordres), le déséquilibre
atteint $\sqrt{\E[T]}$ en moins de $\E[T]$ transactions, l'estimation de $\E[T]$ baisse, et le seuil avec elle~: les \omterm{def:rs:market-data-for-research:bar}{barres}
rétrécissent jusqu'à se clore sur de courtes séquences de flux unilatéral, là où surviennent les grands mouvements du mid.
\end{solution}

\begin{solution}{exo:rs:market-data-for-research:8}
Les deux signaux utilisent des niveaux d'une série ajustée rétroactivement, qui ne sont pas des prix~: chacun a reçu après coup les écarts de roll, de sorte que le rapport de deux
séries ajustées rétroactivement n'est pas le rapport des prix des deux contrats (et un niveau peut être négatif, comme en avril 2020), et la
variation en pourcentage sur douze mois divise une variation en dollars par une base fictive. Le carry doit se calculer à partir des contrats bruts à
chaque date~; la tendance se calcule à partir des rendements ajustés proportionnellement (ou à partir de la variation ajustée rétroactivement divisée par le prix brut détenu).
\end{solution}

\begin{solution}{pb:rs:market-data-for-research:1}
\textbf{1.} 1\,555\,872 messages et 55\,542 transactions~: 28.0 messages par transaction. \textbf{2.} Ajouts 49.3~\%, annulations 47.2~\%,
exécutions 3.6~\%~: presque chaque ordre ajouté est ensuite annulé, car les ordres au marché n'éliminent qu'une faible part du
volume en attente. \textbf{3.} $1\,555\,872 \times 36 = 56.0$ Mo. \textbf{4.} Chaque vue plus grossière peut être reconstruite à partir d'elles, et des \omterm{def:rs:market-data-for-research:bar}{barres} figées
aujourd'hui figeraient une horloge, un seuil et un choix de prix dans toute étude ultérieure. \textbf{5.} Temps 390, ticks 391, volume 388,
montant 388. \textbf{6.} 5.47, 3.58, 4.04 et 3.92. \textbf{7.} 0.94~; $3(1 + 0.94^2) = 5.67$, contre 5.47 mesuré.
\textbf{8.} La volatilité suit le niveau d'activité, et non exactement le volume (et la rafale de nouvelles fait avancer le prix efficient plus vite pour un
volume donné), si bien que les \omterm{def:rs:market-data-for-research:bar}{barres} à volume égal mélangent encore des variances. \textbf{9.} Après l'ouverture, autour de la rafale de nouvelles et avant la clôture,
où le multiplicateur d'activité et la forme en U concentrent les échanges. \textbf{10.} Transactions 0.83, 0.58, 0.96, 0.84~; mid 0.23, 0.39, 0.93,
0.85 (en pour cent au carré). \textbf{11.} Le rebond bid-ask (\cref{prop:rs:market-data-for-research:bounce}). \textbf{12.} Le mid
s'ajuste au prix efficient par paliers, si bien que ses rendements à court horizon sont positivement autocorrélés. \textbf{13.} À partir d'environ 45
secondes (0.95 contre 0.91, puis à moins de 3.4~\%). \textbf{14.} Un estimateur robuste au bruit appliqué aux données à la seconde~: la variance
réalisée à deux échelles, le noyau réalisé ou la prémoyenne. \textbf{15.} \omterm{def:rs:market-data-for-research:bar}{Barres temporelles} de cinq minutes ou plus (ou un estimateur robuste au
bruit) pour la variance de la journée~; une horloge d'événements, ou les messages eux-mêmes, pour les mouvements à court terme. \textbf{16.} Rien ne reflète
l'activité~: elles conservent un calendrier régulier, aligné entre instruments, mais leurs rendements mélangent périodes calmes et périodes actives. \textbf{17.} Un
calendrier commun~: leurs \omterm{def:rs:market-data-for-research:bar}{barres} se terminent à des instants différents selon les instruments et les jours. \textbf{18.} \emph{Résultat
notable~:} les \omterm{def:rs:market-data-for-research:bar}{barres temporelles} d'une minute ont un kurtosis de 5.47, les \omterm{def:rs:market-data-for-research:bar}{barres de montant} 3.92, et à la seconde la variance des prix de transaction vaut
3.6 fois celle du mid. \textbf{19.} Le niveau d'activité aléatoire qui règle à la fois le volume et la vitesse du prix efficient~; avec
\texttt{act\_vol=0} seules la forme en U et la rafale de nouvelles subsistent. \textbf{20.} Aucun~: le prix efficient n'est pas observé, et les transactions,
les mid et les \omterm{def:rs:market-data-for-research:bar}{barres} sont des vues bruitées de celui-ci, chacune avec sa propre erreur.
\end{solution}

\begin{solution}{iq:rs:market-data-for-research:1}
Parce que la variance des prix s'accumule avec les transactions, et non avec le temps~; les \omterm{def:rs:market-data-for-research:bar}{barres de volume} égal (ou de montant égal) ont des rendements bien plus proches
d'une loi gaussienne et moins hétéroscédastiques, ce qui aide les statistiques qui supposent les deux. Le coût~: les \omterm{def:rs:market-data-for-research:bar}{barres} ne s'alignent plus dans le temps
entre instruments, et l'horloge elle-même réagit aux nouvelles.
\iqlookfor{subordination au volume~; la perte d'un calendrier commun.}
\end{solution}

\begin{solution}{iq:rs:market-data-for-research:2}
Bruit de microstructure~: le rebond bid-ask ajoute environ $s^2/2$ par rendement, ce qui domine à haute fréquence. Le signature plot
le met en évidence~; les remèdes consistent à échantillonner moins souvent, à utiliser les mid, ou à recourir à un estimateur robuste au bruit.
\iqlookfor{le rebond, le signature plot, et un remède.}
\end{solution}

\begin{solution}{iq:rs:market-data-for-research:3}
Faire le roll selon une règle qu'une position réelle suivrait (un nombre fixe de jours avant l'échéance, ou sur les positions ouvertes)~; calculer les rendements à partir
du contrat détenu (ajusté proportionnellement), et les signaux sur les niveaux à partir du contrat détenu à cette date. Les niveaux ajustés rétroactivement réécrivent le
passé à chaque roll et peuvent devenir négatifs~; leurs variations en pourcentage sont fausses.
\iqlookfor{rendements contre niveaux, et le fait que l'ajustement est un choix de modélisation.}
\end{solution}

\begin{solution}{iq:rs:market-data-for-research:4}
Messages bruts immuables, avec horodatages de la bourse et de réception et numéros de séquence, partitionnés par date et par instrument,
compressés~; environ 76 To par an à l'état brut, à 40 octets par message et 1.5 million de messages par instrument et par jour. Tables dérivées (carnets, \omterm{def:rs:market-data-for-research:bar}{barres}) reconstruites par
code versionné, stockées en colonnes, partitionnées selon leur mode de lecture, chacune avec sa filiation enregistrée.
\iqlookfor{immutabilité, filiation, le partitionnement, et un ordre de grandeur de la taille.}
\end{solution}

\begin{solution}{iq:rs:market-data-for-research:5}
Clôtures non synchrones~: la dernière transaction d'un titre peu liquide peut dater de plusieurs minutes, ce qui crée de faux décalages lead--lag et réduit
les corrélations (effet Epps)~; le rebond des prix de clôture~; l'hétéroscédasticité due à une activité inégale~; les transactions d'enchères dans les premières
et dernières \omterm{def:rs:market-data-for-research:bar}{barres}.
\iqlookfor{non-synchronicité et bruit, et pas seulement le look-ahead.}
\end{solution}

\begin{solution}{iq:rs:market-data-for-research:6}
Sachant $V$, on a $r \sim \mathcal N(0, \sigma^2V)$, donc $\E[r^2] = \sigma^2\E[V]$ et $\E[r^4] = 3\sigma^4\E[V^2]$~; le rapport vaut $3\E[V^2]/\E[V]^2
= 3(1 + \mathrm{CV}^2)$.
\iqlookfor{conditionnement sur le volume et moment d'ordre quatre d'une loi normale.}
\end{solution}
